\begin{abstract}
We introduce \textsc{QLoRA}, an efficient finetuning methodology that dramatically lowers memory requirements, enabling the finetuning of a 65B parameter model on a single 48GB GPU while maintaining performance on par with conventional 16-bit finetuning. \textsc{QLoRA} operates by propagating gradients through a static, 4-bit quantized version of a pretrained language model and optimizing Low Rank Adapters~(LoRA). Our leading suite of models, named \textbf{Guanaco}, surpasses all prior open-access models on the Vicuna benchmark, achieving 99.3\% of ChatGPT's performance after just 24 hours of single-GPU finetuning. To achieve these results, \textsc{QLoRA} introduces several key memory-saving techniques without compromising effectiveness: (a) the 4-bit NormalFloat (NF4) datatype, which is theoretically optimal for parameters with a normal distribution, (b) Double Quantization, a process that further shrinks memory consumption by quantizing the quantization parameters themselves, and (c) Paged Optimizers, which control memory usage surges. Leveraging \textsc{QLoRA}, we finetuned over 1,000 models, conducting an extensive study of instruction following and chatbot quality across 8 instruction datasets, diverse architectures (including LLaMA and T5), and model sizes previously impractical for finetuning (notably 33B and 65B parameters). Our findings demonstrate that \textsc{QLoRA} achieves state-of-the-art results using modestly sized, high-quality datasets—even with models smaller than former top performers. In-depth evaluations, utilizing both human raters and GPT-4, indicate that GPT-4 can offer a cost-effective and reliable alternative to manual assessments. Moreover, we observe that popular chatbot benchmarks fail to accurately reflect real chatbot capabilities. An illustrative “lemon-picked” analysis identifies specific shortcomings of \textbf{Guanaco} relative to ChatGPT. We openly provide all our models, source code, and custom CUDA kernels supporting 4-bit training.\footnote{\url{https://github.com/artidoro/qlora} and \url{https://github.com/TimDettmers/bitsandbytes}}
\end{abstract}

\section{Introduction}

Adapting large language models (LLMs) through finetuning is one of the most effective approaches to enhancing their capabilities \citep{min2021metaicl, wei2021finetuned, ouyang2022training, wang2022super, wang2022self, liu2022few} and for instilling or suppressing specific behaviors \citep{ouyang2022training,askell2021general,bai2022training}. Nonetheless, the costs associated with finetuning extremely large models are immense; for example, standard 16-bit finetuning for a LLaMA model with 65B parameters~\citep{touvron2023llama} consumes over 780 GB of GPU memory. Though advances in quantization offer notable memory savings at inference \citep{dettmers2022llm, dettmers2022case, frantar2022gptq, xiao2022smoothquant}, such strategies fail to maintain effectiveness during training and typically become unstable \citep{wortsman2023stable}.

In this work, we show for the first time that a 4-bit quantized model can be finetuned without sacrificing any performance. The proposed method, \textsc{QLoRA}, introduces a new, high-precision quantization strategy for pretrained models to the 4-bit level, followed by the integration of compact, trainable Low-rank Adapter weights \citep{hu2021lora} 

which are updated by propagating gradients through the quantized parameter space.

\textsc{QLoRA} achieves a drastic reduction in the memory demand for finetuning a 65B parameter model, lowering the requirement from over 780GB to under 48GB of GPU memory, all while sustaining both computational efficiency and predictive accuracy on par with conventional 16-bit finetuning. This development significantly broadens the accessibility of LLM finetuning, enabling the largest open models to be finetuned on a single GPU. Leveraging \textsc{QLoRA}, we developed the \textbf{Guanaco} model series; notably, our second most capable model attains 97.8\% of ChatGPT's performance according to the Vicuna~\citep{vicuna2023} benchmark, completing training in under 12 hours on a consumer-grade GPU. Using a professional GPU for 24 hours, our largest model reaches 99.3\% performance, effectively matching ChatGPT's results on the same benchmark. When deployed, our smallest \textbf{Guanaco} variant (7B parameters) operates with only 5 GB of memory and surpasses the 26 GB Alpaca model by more than 20 percentage points on the Vicuna benchmark (see Table~\ref{tbl:vicuna_eval}).

\textsc{QLoRA} presents several advancements aimed at decreasing memory consumption while maintaining high model performance: (1) \textbf{4-bit NormalFloat}, an information-theoretically optimal quantization format tailored for data with a normal distribution, which outperforms both 4-bit Integer and 4-bit Float representations in practice. (2) \textbf{Double Quantization}, which involves further quantizing the quantization coefficients themselves, resulting in average savings of roughly 0.37 bits per parameter (translating to about 3 GB for a 65B-parameter model). (3) \textbf{Paged Optimizers}, leveraging NVIDIA's unified memory to mitigate memory surges typically seen with gradient checkpointing during training with long sequence lengths. Collectively, these features are integrated within a refined LoRA framework that incorporates adapters at all layers of the network, effectively mitigating most of the accuracy compromises observed in earlier studies.

The memory efficiency provided by \textsc{QLoRA} allows for a comprehensive exploration of instruction finetuning and chatbot effectiveness at scales unmanageable with conventional finetuning approaches due to memory constraints. We thus train over 1,000 distinct models, spanning various instruction-tuning datasets, architectures, and parameter scales from 80M to 65B. Beyond demonstrating that \textsc{QLoRA} matches the performance of standard 16-bit training (\S\ref{sec:comp_to_std}) and successfully trains the leading chatbot \textbf{Guanaco} (\S\ref{sec:pushing}), we further investigate patterns arising in the trained models. Our findings reveal that dataset quality supersedes sheer quantity: for example, a 9k-example dataset (OASST1) surpassed a 450k-example dataset (FLAN v2, subsampled) in terms of chatbot capability, even when both datasets aimed to foster generalization in instruction following. Additionally, our analysis shows that excelling on the Massive Multitask Language Understanding (MMLU) benchmark does not necessarily equate to superior performance on the Vicuna chatbot benchmark, highlighting that appropriateness of dataset selection for the specific application is more crucial than dataset size.

In addition, we conduct a comprehensive evaluation of chatbot systems leveraging assessments from both human reviewers and GPT-4. We adopt a tournament-style approach for benchmarking, in which models are pitted against each other to generate optimal responses for specified prompts. The determination of match winners is carried out either by GPT-4 or by human judges. These tournament match outcomes are synthesized into Elo scores~\citep{elo1967proposed, elo1978rating}, which serve as a metric for ranking chatbot capabilities. Our findings reveal substantial concordance between GPT-4 and human raters concerning model rankings within tournaments; nevertheless, we also observe notable cases of divergence. This underscores that, while automated model evaluations offer a cost-effective alternative to manual annotation, they introduce a degree of uncertainty.

To supplement our quantitative benchmarks, we include a qualitative investigation of the \textbf{Guanaco} models. Through this analysis, we identify examples of both effective and unsuccessful model behavior that are not reflected by numerical benchmarking alone.

For the benefit of future research, we publicly share all generated outputs, along with accompanying human and GPT-4 evaluations. Our codebase and CUDA kernels are made available as open-source, with integration into the Hugging Face transformers framework \citep{wolf2019huggingface} to enhance accessibility. We also distribute adapter collections for models with 7/13/33/65B parameters, trained across 8 distinct instruction-following datasets, yielding 32 openly released, fine-tuned model variants.

\section{Background}
\paragraph{Block-wise k-bit Quantization} The procedure of quantization involves transforming data from a high-precision representation to one of reduced precision, commonly by reducing the number of bits per element—for example, converting 32-bit floating point values to 8-bit integers. To guarantee effective usage of the reduced precision's dynamic range, it is standard practice to rescale the input—typically organized as tensors—so that the original values are normalized relative to their maximum absolute value, thereby aligning them to the range of the target data type. For instance, converting an FP32 tensor into an Int8 tensor with values in $[-127, 127]$ is accomplished as follows:

\begin{equation}
    \mathbf{X}^{ \text{Int8}} = \text{round}\left( \frac{127}{{\text{absmax}}(\mathbf{X}^{{\text{FP32}}})}\mathbf{X}^{{\text{FP32}}} \right) = \text{round}(c^{\text{FP32}}\cdot \mathbf{X}^{{\text{FP32}}}),
\end{equation}

where the parameter $c$ represents the {\em quantization constant} or {\em quantization scale}. The original values can be retrieved via the inverse operation, dequantization:

\begin{equation}
    \text{dequant}(c^{\text{FP32}}, \mathbf{X}^{\text{Int8}}) = \frac{\mathbf{X}^{\text{Int8}}}{c^{\text{FP32}}} = \mathbf{X}^{\text{FP32}}
\end{equation}

One limitation of this method is its inefficiency in utilizing quantization bins when the input tensor contains values of exceptionally high magnitude (i.e., outliers). In such situations, several bit patterns—corresponding to specific quantization bins—may end up with few or even no input values mapped to them. To alleviate problems caused by these outliers, a prevalent strategy involves dividing the input tensor into smaller segments, or blocks, each subjected to quantization separately with its own quantization constant $c$. More precisely, the input tensor $\mathbf{X}\in \mathbb{R}^{b\times h}$ is reshaped into $n$ consecutive blocks, each of size $B$, by first flattening $\mathbf{X}$ and partitioning the resulting array into ${n = ({b\times h} )/{B} }$ distinct blocks. Quantization is then carried out independently for each block, using Equation~1, resulting in a quantized tensor together with $n$ quantization constants, one for each block, denoted as $c_i$.

\paragraph{Low-rank Adapters} The LoRA finetuning technique \citep{hu2021lora} seeks to minimize memory overhead during adaptation by introducing a compact subset of trainable parameters, known as adapters, while the majority of the model’s parameters remain unaltered. During the course of stochastic gradient descent, the gradients propagate through the unchanged, pretrained model weights, focusing the updates exclusively on the adapters so as to refine the model with respect to the loss objective. The LoRA approach supplements a standard linear projection with an additional projection decomposed into low-rank factors. Specifically, for a projection ${\mathbf{X} \mathbf{W} = \mathbf{Y}}$ where $\mathbf{X}\in  \mathbb{R}^{b\times h}$ and $\mathbf{W}\in  \mathbb{R}^{h\times o}$, the computation in LoRA is given by:

\begin{equation}
    \mathbf{Y} = \mathbf{X}\mathbf{W} + s\mathbf{X}\mathbf{L}_1\mathbf{L}_2,
\end{equation}

with $\mathbf{L}_1\in  \mathbb{R}^{h\times r}$, $\mathbf{L}_2\in  \mathbb{R}^{r\times o}$, and $s$ representing a scaling constant.

\paragraph{Memory Requirement of Parameter-Efficient Finetuning} A key aspect to consider is the memory consumption associated with LoRA during training, particularly with respect to both the count and dimensionality of adapters employed. Owing to LoRA’s extremely lightweight memory demands, it is feasible to deploy multiple adapters for enhanced performance while maintaining a relatively stable overall memory usage. Although LoRA is conceptualized as a Parameter Efficient Finetuning (PEFT) technique, the predominant share of memory usage in LLM finetuning arises from storing activation gradients rather than the LoRA parameters themselves. For example, in a scenario where a 7B LLaMA model is trained on FLAN v2 with a batch size of 1, and the LoRA weights represent the widely adopted 0.2\% proportion of the full model weights \citep{hu2021lora,liu2022few}, the input gradients attributable to LoRA account for 567 MB of memory, whereas the parameters themselves only require 26 MB. Employing gradient checkpointing~\citep{chen2016training} can decrease the input gradient storage to an average of 18 MB per sequence, making this component still larger than the aggregate memory cost of the LoRA weights. In contrast, the base model with 4-bit precision occupies 5,048 MB. These figures underscore the importance of gradient checkpointing and suggest that further minimizing LoRA parameter size yields limited gains in memory efficiency. Consequently, the addition of further adapters has minimal impact on total training memory needs (a detailed analysis is provided in Appendix~\ref{app:memory}). As elaborated later, this flexibility is essential for attaining performance equivalent to full 16-bit precision.

\section{\textsc{QLoRA} Finetuning}

\textsc{QLoRA} enables high-fidelity finetuning using 4-bit representations through the use of two principal innovations: 4-bit NormalFloat (NF4) quantization and Double Quantization. Alongside these, we present Paged Optimizers as a strategy to mitigate memory overflow during gradient checkpointing, an issue which has historically limited the feasibility of large model finetuning on a single node due to memory constraints.

In the \textsc{QLoRA} framework, model weights are maintained in a low-precision storage format, typically 4 bits, while computations are carried out in a higher-precision format, usually BFloat16. Operationally, this involves dequantizing \textsc{QLoRA} weight tensors to BFloat16 whenever computation is required, followed by executing the associated matrix operations in 16-bit precision.

The subsequent sections elaborate on the technical components that comprise \textsc{QLoRA}, leading up to its formal specification.

\paragraph{4-bit NormalFloat Quantization} The NormalFloat (NF) type expands upon Quantile Quantization \citep{dettmers2022optimizers}, an information-theoretically ideal approach in which each bin in the quantized space is allocated the same number of entries from the input tensor. Quantile quantization relies on calculating the empirical cumulative distribution function to determine the quantiles of the input tensor.

A significant drawback of quantile quantization is the computational cost associated with accurate quantile estimation. To address this, accelerated quantile approximation schemes such as SRAM quantiles~\citep{dettmers2022optimizers} are often utilized. However, these approximations introduce considerable quantization error, particularly for outlier values—elements which frequently hold substantial importance.

Both the computational overhead and the risk of large approximation errors are circumvented when input tensors are sampled from distributions that are consistent modulo a quantization constant; in these scenarios, tensors share identical quantiles, making the use of precise quantile estimates tractable.

Pretrained neural network weights are typically distributed according to a zero-centered normal distribution with standard deviation $\sigma$ (refer to Appendix~\ref{app:normality}). To map all weights into a uniform distribution, we can adjust $\sigma$ so that the entire distribution aligns perfectly within the numerical limits of our chosen data type. In our implementation, we adopt $[-1, 1]$ as this target range. Consequently, normalization of both the neural network weights and the data type's quantiles into the $[-1, 1]$ interval becomes necessary.

For zero-mean normal distributions with arbitrary $\sigma$ constrained to $[-1, 1]$, the theoretically most efficient data type from an information-theoretic perspective can be constructed in the following way: (1) identify the $2^k+1$ quantile boundaries for the standard normal distribution $N(0, 1)$, thereby defining a quantile-based quantization for $k$-bit precision; (2) rescale these quantile points so that their values reside within $[-1, 1]$; (3) normalize any given input weight tensor via absolute maximum scaling so that its values, too, are within $[-1, 1]$, before quantization proceeds.

With the ranges of both the data type and the weights synchronized, traditional quantization can then be applied. The rescaling in step (3) effectively aligns the standard deviation of the tensor with that of the $k$-bit quantized representation. To formalize this, the $2^k$ representative values $q_i$ for the data type are computed by

\begin{equation}
q_i  = \frac{1}{2}\left( Q_X\left(\frac{i}{2^k+1}\right) + Q_X\left(\frac{i+1}{2^k+1}\right)\right),
\end{equation}

where $Q_X(\cdot)$ denotes the quantile function associated with the standard normal distribution $N(0,1)$. One limitation of symmetric k-bit quantization schemes is the lack of a precise representation of zero, a critical attribute when exact quantization of padding and other zero-valued components is required. To address this and simultaneously leverage all $2^k$ quantization levels for a k-bit format, we design an asymmetric type by deriving quantiles $q_i$ across two intervals: $2^{k-1}$ quantiles for the negative range and $2^{k-1}+1$ quantiles for the positive range. These two quantile sets are merged, and one of the overlapping zeros is eliminated. We refer to this newly constructed type, which ensures an equal expected allocation of values in each bin, as \emph{k-bit NormalFloat} (NFk), reflecting its optimality (in the information-theoretic sense) for data with a zero-centered normal distribution. Further details about the precise value set can be found in Appendix~\ref{app:nf4}.

\paragraph{Double Quantization{}}
We propose \emph{Double Quantization} (DQ) as a strategy to quantize the quantization constants themselves, thus further decreasing memory requirements. Achieving high fidelity with 4-bit quantization generally necessitates small blocksizes~\citep{dettmers2022case}, but this comes at the cost of increased memory for storing the quantization constants. For example, if 32-bit constants are employed with a blocksize of 64 for $\mathbf{W}$, the additional memory due to constants averages $32/64 = 0.5$ bits per parameter. Double Quantization significantly mitigates this memory overhead.

The method operates as follows: quantization constants from the initial quantization, $c_2^{\text{FP32}}$, are taken as input for a subsequent quantization stage, resulting in quantized constants $c_2^{\text{FP8}}$ and a new set of quantization constants $c_1^{\text{FP32}}$. For this second quantization step, 8-bit floating point numbers with a blocksize of 256 are utilized, as no loss in accuracy is detected under 8-bit quantization, echoing findings from \citet{dettmers2022case}. Since all $c_2^{\text{FP32}}$ values are non-negative, their mean is subtracted before the second quantization step to center the distribution around zero, enabling symmetric quantization. In this setup with blocksize 64, the overall memory consumption per parameter drops from $32/64=0.5$ bits to $8/64 + 32/(64\cdot256) = 0.127$ bits, achieving a savings of 0.373 bits per parameter.

\paragraph{Paged Optimizers} leverage NVIDIA's unified memory~\footnote{\tiny \url{https://docs.nvidia.com/cuda/cuda-c-programming-guide}}, a capability that seamlessly manages page transfers between CPU and GPU memory, ensuring reliable GPU operations even when GPU memory becomes insufficient. This mechanism resembles conventional paging that occurs between CPU RAM and storage drives. In our approach, optimizer state memory is allocated as paged memory, allowing it to be transparently moved from the GPU to the CPU RAM when necessary and subsequently reloaded onto the GPU during optimizer updates when the data is required.

\paragraph{\textsc{QLoRA}.} Building on the previously outlined elements, we define \textsc{QLoRA} applied to a single linear transformation in the quantized foundational model, enhanced by a single LoRA adapter, as:

\begin{equation}
    \mathbf{Y}^{\text{BF16}} =  \mathbf{X}^{\text{BF16}} \text{doubleDequant}(c_1^{\text{FP32}}, c_2^{\text{k-bit}}, \mathbf{W}^{\text{NF4}}) + \mathbf{X}^{\text{BF16}}\mathbf{L}^{\text{BF16}}_1\mathbf{L}^{\text{BF16}}_2,
\end{equation}

The doubleDequant$(\cdot)$ function is specified by:

\begin{equation}
    \text{doubleDequant}(c_1^{\text{FP32}}, c_2^{\text{k-bit}}, \mathbf{W}^{\text{k-bit}}) = \text{dequant}(\text{dequant}(c_1^{\text{FP32}}, c_2^{\text{k-bit}}), \mathbf{W}^{\text{4bit}}) = \mathbf{W}^{\text{BF16}},
\end{equation}

In our configuration, $\mathbf{W}$ is represented in NF4, while $c_2$ is stored using the FP8 format.

To improve quantization fidelity, we select a blocksize of 64 for $\mathbf{W}$. For $c_2$, a larger blocksize of 256 is adopted to reduce memory usage.

During parameter optimization, it suffices to compute the gradients $\frac{\partial E}{\partial \mathbf{L}_i}$ pertaining to the LoRA adapter weights, omitting the gradients for the 4-bit weights $\frac{\partial E}{\partial \mathbf{W}}$. Nonetheless, computing $\frac{\partial E}{\partial \mathbf{L}_i}$ requires intermediate computation of $\frac{\partial \mathbf{X}}{\partial \mathbf{W}}$, which is executed via equation~(5). This step involves dequantizing from the stored representation $\mathbf{W}^{\text{NF4}}$ to the computational data type $\mathbf{W}^{\text{BF16}}$, allowing the calculation of $\frac{\partial \mathbf{X}}{\partial \mathbf{W}}$ at BFloat16 precision.

In summary, \textsc{QLoRA} utilizes a single storage data format—most commonly 4-bit NormalFloat—and a separate computation data type, specifically 16-bit BrainFloat. To execute both the forward and backward passes, the stored weights are dequantized from the storage format to the computation format. However, during gradient computation, only the LoRA parameters, maintained in 16-bit BrainFloat, are updated.

\section{QLoRA vs. Standard Finetuning}
\label{sec:comp_to_std}

After outlining the mechanism of QLoRA and its capacity for significant memory savings during model finetuning, an essential question remains: can QLoRA match the effectiveness of conventional full-model finetuning? Additionally, it is important to dissect QLoRA's components, such as the role of NormalFloat4 relative to standard Float4 representations. The next sections describe the experiments we conducted to address these inquiries.

\paragraph{Experimental setup.} We examine three neural network architectures—encoder, encoder-decoder, and decoder-only—contrasting QLoRA with both 16-bit adapter-based finetuning and complete finetuning on models as large as 3B parameters. The tasks in our evaluation comprise GLUE \citep{wang2018glue} with RoBERTa-large \citep{liu2019roberta}, Super-NaturalInstructions (TKInstruct) \citep{wang2022super} leveraging T5 \citep{t5}, and 5-shot MMLU \citep{hendrycksmeasuring} following LLaMA finetuning on Flan v2 \citep{longpre2023flan} and Alpaca \citep{alpaca}. To further investigate the relative advantages of NF4 against other 4-bit quantization schemes, we adopt the experimental framework of \citet{dettmers2022case} to assess post-quantization zero-shot accuracy and perplexity over a range of model families (OPT~\citep{zhang2022opt}, LLaMA~\citep{touvron2023llama}, BLOOM~\citep{scao2022bloom}, Pythia~\citep{biderman2023pythia}) covering model sizes from 125m to 13B.

Additional details for each experiment are given in the results section for greater clarity, with complete information provided in Appendix~\ref{app:exp-setup}.

Although paged optimizers are crucial for training 33B and 65B \textsc{QLoRA} models on single GPUs with 24GB or 48GB memory, we omit detailed measurements of Paged Optimizer performance due to the infrequency of paging, which only happens with mini-batches that have long sequence lengths. Nevertheless, we report a runtime assessment for paged optimizers with 65B models on 48GB GPUs, observing that, with a batch size of 16, their training speed is on par with that of standard optimizers. More systematic studies are needed in the future to determine when and why paging-induced slowdowns may occur.

\paragraph{Default LoRA hyperparameters do not match 16-bit performance} 

Employing the typical approach of applying LoRA exclusively to the query and value attention projection matrices \citep{hu2021lora} does not yield full finetuning-level results for large-scale base models. As illustrated in Figure~\ref{fig:hyper} for the finetuning of LLaMA 7B on Alpaca, our findings indicate that the number of LoRA adapters deployed—specifically, ensuring that adapters are present in every linear transformer block layer—is pivotal to reproducing full finetuning performance. Other LoRA settings, such as the projection dimension $r$, have negligible effect on outcomes (see Appendix \ref{app:exp-setup} for further detail).

In a similar vein, our analysis indicates that the standard hyperparameters used for fully finetuned baseline models are insufficiently optimized. To establish solid baselines, we conduct a comprehensive hyperparameter sweep, varying the learning rate between 1e-6 and 5e-5, and experimenting with batch sizes ranging from 8 to 128. Figure~\ref{fig:hyper} presents these findings for 7B LLaMA finetuning on Alpaca.

\paragraph{4-bit NormalFloat Outperforms 4-bit Floating Point}

Although the 4-bit NormalFloat (NF4) format is theoretically optimal in terms of information retention, its practical benefits require empirical validation. Adopting the methodology from \citet{dettmers2022case}, we evaluate quantized LLMs—including OPT \citep{zhang2022opt}, BLOOM \cite{scao2022bloom}, Pythia \cite{biderman2023pythia}, and LLaMA—spanning sizes from 125M to 65B parameters and utilizing various data types. Performance on language modeling and multiple zero-shot tasks is reported in Figure~\ref{fig:nf4} and Table~\ref{tbl:nf4_ppl}. The results demonstrate that NF4 markedly surpasses FP4 and Int4 in performance, and that double quantization effectively decreases memory requirements while preserving model accuracy.

\paragraph{k-bit \textsc{QLoRA} Achieves Parity with 16-bit Full Finetuning and 16-bit LoRA}

Prior studies have shown that although 4-bit quantization enables \emph{inference}, it typically comes at the cost of diminished performance compared to 16-bit models \citep{dettmers2022case,frantar2022gptq}. This observation raises the important question of whether 4-bit adapter finetuning can restore any lost accuracy. To explore this, we perform two sets of experiments.

The first involves benchmarking full 16-bit finetuning of RoBERTA and T5 models, ranging from 125M to 3B parameters, on both the GLUE and Super-NaturalInstructions datasets. Table~\ref{tab:nf4vsbf16} summarizes the outcomes. In both benchmark suites, we find that adapter-based approaches at 16-bit, 8-bit, and 4-bit precision consistently match the performance of fully finetuned 16-bit models. These findings suggest that adapter finetuning, even after aggressive quantization, is able to fully recoup performance deficits introduced by lower precision.

In our second experimental configuration, because full finetuning of models with 11B parameters or more exceeds the memory capacity of a single GPU server, we further investigate whether 4-bit \textsc{QLoRA} can achieve parity with 16-bit LoRA when scaling model sizes from 7B to 65B parameters. For this purpose, we perform finetuning on LLaMA models (ranging from 7B to 65B parameters) using two instruction-tuned datasets, Alpaca and FLAN v2, and assess performance on the MMLU benchmark employing 5-shot accuracy. Table~\ref{tbl:llama-datatype} presents these outcomes, demonstrating that NF4 with double quantization effectively restores the MMLU performance attained by 16-bit LoRA. We also observe that \textsc{QLoRA} configured with FP4 trails the 16-bit brain float LoRA baseline by roughly 1 percentage point. These results substantiate our two main conclusions: (1) \textsc{QLoRA} using NF4 is capable of matching both 16-bit full finetuning and 16-bit LoRA finetuning results, and (2) NF4 outperforms FP4 in terms of quantization accuracy.

\paragraph{Summary} Collectively, our findings indicate that utilizing 4-bit \textsc{QLoRA} with the NF4 datatype can deliver equivalent results to 16-bit full and 16-bit LoRA finetuning on widely recognized academic benchmarks. We further demonstrate that NF4 yields higher accuracy than FP4, and that double quantization introduces no significant loss in model performance. These points collectively provide strong evidence that 4-bit \textsc{QLoRA} is a robust alternative to 16-bit approaches.

Consistent with existing literature on quantization \citep{dettmers2022case}, the results from both our MMLU and Elo evaluations suggest that, given a fixed finetuning and inference resource budget, increasing model parameter count while employing reduced precision can be advantageous. This underscores the practical efficiency advantages provided by \textsc{QLoRA}. Since our 4-bit finetuning experiments do not reveal any noticeable drop in accuracy compared to full finetuning, it remains an open question as to the precise location of the performance-precision trade-off for QLoRA tuning—a topic we defer for future investigation.

Our research advances the study of instruction tuning by extending it to scales unattainable with conventional 16-bit finetuning on standard academic hardware.

\section{Pushing the Chatbot State-of-the-art with QLoRA}
\label{sec:pushing}
After demonstrating that 4-bit \textsc{QLoRA} attains performance parity with 16-bit approaches across model sizes, task domains, and benchmarks, we delve deeper into the instruction finetuning of the largest open-source language models available to researchers. To comprehensively evaluate instruction tuning, we leverage the demanding Natural Language Understanding benchmark (MMLU) and introduce new strategies for assessing chatbot efficacy in real-world settings.

\subsection{Experimental setup} Here, we outline the essential elements of our experimental procedure; additional methodological specifics are presented in Appendix \ref{app:chatbot}.

\paragraph{Data} Since there appears to be a gap in systematic analyses of recent instruction-following corpora, we compile a selection of eight contemporary datasets. These encompass datasets created via crowd-sourcing (OASST1~\citep{kopf2023openassistant}, HH-RLHF~\citep{bai2022training}), those resulting from model distillation (Alpaca~\citep{alpaca}, self-instruct~\citep{wang2022self}, unnatural-instructions~\citep{honovich2022unnatural}), collections assembled from diverse sources (FLAN v2~\citep{chung2022scaling}), as well as composite datasets (Chip2~\citep{laion2023}, Longform~\citep{koksal2023longform}). This suite captures a variety of languages, dataset magnitudes, and licensing schemes.

\paragraph{Training Setup} To ensure that variation in training objectives does not confound results, we standardize finetuning across all datasets using QLoRA with cross-entropy loss under supervised learning, deliberately omitting reinforcement learning—even when human rating data is present. Where datasets delineate instruction from response, we focus finetuning solely on responses (see our ablation analyses in Appendix \ref{app:chatbot}). OASST1 and HH-RLHF present multiple response options; in these cases, we choose the highest-ranked reply per node in the conversation hierarchy and train on the aggregated conversation path, encompassing the associated instructions. Throughout, we employ the NF4 variant of \textsc{QLoRA}, utilizing both double quantization and paged optimizers to mitigate memory surges during gradient checkpointing. Limited hyperparameter exploration is conducted for LLaMA models at 13B and 33B scales, revealing that parameters tuned at 7B transfer broadly (including epoch count), except for batch size and learning rate. Specifically, for 33B and 65B, we use half the learning rate and twice the batch size as compared to 7B.

\paragraph{Baselines} Our evaluations include comparisons with both non-commercial models (Vicuna~\citep{vicuna2023}, Open Assistant~\citep{kopf2023openassistant}) and proprietary systems (GPT-4~\citep{openai2023gpt}, GPT-3.5-turbo, Bard). The Open Assistant reference point is based on the LLaMA 33B architecture, finetuned using Reinforcement Learning from Human Feedback (RLHF) on OASST1—the same corpus as ours. Meanwhile, Vicuna utilizes full fine-tuning of LLaMA 13B on user-generated content from ShareGPT, constituting a distilled variant of the OpenAI GPT series.

\subsection{Evaluation}

Consistent with established protocols, we adopt the MMLU (Massively Multitask Language Understanding) benchmark~\citep{hendrycksmeasuring} to gauge linguistic comprehension across diverse domains. This suite presents multiple-choice problems spanning 57 disciplines, such as elementary mathematics, U.S. history, computer science, and law. Results are reported as 5-shot test accuracies.

Our evaluation of generative language proficiency encompasses both automated and human assessments. This latter suite of evaluations utilizes human-crafted queries and is intended to gauge the quality of model outputs. Although increasingly regarded as a realistic means of testing chatbot model efficacy, the field lacks a universally established methodology. Below, we outline our chosen evaluation approach, consistently employing nucleus sampling with $p=0.9$ and a temperature setting of $0.7$.

\paragraph{Benchmark Data} We perform our evaluations using two specifically assembled question datasets: the Vicuna prompts \citep{vicuna2023} and the OASST1 validation dataset \citep{kopf2023openassistant}. For the Vicuna prompts, which comprise 80 questions from a broad array of topics, we use the dataset in its original form. The OASST1 dataset features a multilingual, crowdsourced collection of multi-turn exchanges between users and assistants. Here, every user message from the validation set serves as a distinct query, with preceding conversational turns incorporated into the prompt, resulting in 953 unique queries. We refer to these datasets as the Vicuna and OA benchmarks, respectively.

\paragraph{Automated Evaluation} Initially, following the assessment methodology proposed by \citet{vicuna2023}, we leverage GPT-4 to evaluate different systems relative to ChatGPT (GPT-3.5 Turbo) on the Vicuna benchmark. For each query, GPT-4 receives both ChatGPT's and the target model's replies and assigns each a score out of ten, accompanied by a rationale. The model's aggregate performance is represented as a percentage of ChatGPT’s score. It is possible for this relative score to surpass 100\% should the evaluated model obtain a higher raw score than ChatGPT. We observe a notable bias wherein GPT-4 tends to favor the response positioned first in the prompt. To mitigate such positional biases, we advise averaging results over both possible presentation orders.

Subsequently, we assess models through direct output comparisons. Here, the evaluation is recast as a three-category classification problem to accommodate ties. GPT-4 is tasked with selecting the superior response or indicating a tie, supported by an explanation. These direct, pairwise assessments are performed exhaustively across all system pairs within both the Vicuna and OA benchmarks.

\paragraph{Human Evaluation} Despite emerging studies suggesting that generative models can function as effective evaluators for system performance \citep{fu2023gptscore}, there is, to date, insufficient evidence that GPT-4’s ratings are congruent with human judgments when assessing chatbot performance. As a result, we carry out two parallel human evaluation protocols on the Vicuna benchmark, each mirroring one of the automated approaches outlined above. For these tasks, we employ Amazon Mechanical Turk (AMT), assigning two annotators per comparison with ChatGPT and three annotators per pairwise system comparison.

\paragraph{Elo Rating} We organize both human and automated pairwise model assessments into a tournament framework, wherein models are pitted against each other in head-to-head matches. During each match, two models are evaluated on their ability to generate the superior response to a shared prompt. This methodology aligns with the approaches employed by \citet{bai2022training} and \citet{vicuna2023}, but we extend it by including GPT-4-based evaluations alongside human judgments. To calculate the Elo~\citep{elo1967proposed,elo1978rating} scores, we draw random samples from the pool of annotated pairwise comparisons. The Elo system, commonly applied in chess and various competitive games, estimates a player's expected success rate relative to another's. For instance, if a model holds an Elo of 1100 against an opponent at 1000, it is predicted to win about 65\% of the time; conversely, two models with equal scores (e.g., 1000 vs 1000 or 1100 vs 1100) are each anticipated to win half the matches. Elo ratings are updated after each match according to the deviation from the anticipated outcome: surprising upsets yield substantial Elo adjustments, while expected results lead to more modest changes. As matches accumulate, the Elo rating converges toward accurately reflecting each model’s proficiency in this task. All models are initialized with a rating of 1,000 and we adopt $K=32$ for the update parameter. Following the practice of \citet{vicuna2023}, this evaluation is iterated 10,000 times using different random seeds to minimize the influence of sequence effects, such as the initial pairing of specific models.

\subsection{\textbf{Guanaco}: \textsc{QLoRA} trained on OASST1 Represents the Cutting Edge in Chatbots}

Through a combination of both automated and manual assessments, our analysis reveals that the leading \textsc{QLoRA}-fine-tuned model, {Guanaco} 65B—finetuned on a modified OASST1 dataset—emerges as the highest-performing open-source chatbot, delivering results on par with ChatGPT. In head-to-head comparisons with GPT-4, human annotators assign {Guanaco} 65B and 33B an average win probability of 30\% as determined by Elo scores, which marks the best performance reported so far for open-source alternatives.

Table~\ref{tbl:vicuna_eval} presents the outcomes on the Vicuna benchmark \cite{vicuna2023} relative to ChatGPT. Notably, {Guanaco} 65B achieves the second-highest performance after GPT-4, registering a score at 99.3\% of ChatGPT’s. While {Guanaco} 33B utilizes more parameters than Vicuna 13B, it leverages 4-bit quantization, resulting in superior memory efficiency (21 GB compared to 26 GB) while improving performance over Vicuna 13B by three percentage points. Additionally, {Guanaco} 7B is compact enough to run on current smartphones (5 GB memory footprint), yet outperforms Alpaca 13B by nearly 20 percentage points.

Nevertheless, the confidence intervals in Table~\ref{tbl:vicuna_eval} are substantial, leading to significant overlap among models. We suggest this variance is largely attributable to ambiguities in scale definition—for instance, the interpretation of an “8” on a 10-point scale can vary across contexts. To mitigate such inconsistencies, we advocate for the Elo rating system \citep{elo1967proposed}, which relies on \textit{pairwise} judgments by human annotators and GPT-4, bypassing the pitfalls of using an absolute scale. Table~\ref{tbl:elo} displays Elo scores for the top models. There is some discrepancy between human and GPT-4 rankings, most notably for Guanaco 7B, but overall concordance is observed across most models, as indicated by a system-level Kendall Tau of $\tau=0.43$ and a Spearman correlation of $r=0.55$. On the example level, agreement drops, with a Fleiss $\kappa=0.25$ for majority voting between human annotators and GPT-4. In summary, these findings suggest moderate alignment in system-level evaluations by GPT-4 and humans, supporting the utility of model-driven evaluation as a reasonably dependable alternative to manual assessments. Further discussion can be found in Section~\ref{ssec:considerations}.

The Elo rankings displayed in Table~\ref{tbl:elo_large} reveal that {Guanaco} models at the 33B and 65B scales surpass all models except for GPT-4 on both the Vicuna and OA benchmarks, and their performance is on par with ChatGPT, as also reflected in Table~\ref{tbl:vicuna_eval}. It is important to highlight that the Vicuna benchmark tends to favor open-source models, while ChatGPT has an advantage in the broader OA benchmark. Additionally, analysis of Tables \ref{tbl:mmlu-llama} and \ref{tbl:vicuna_eval} demonstrates that the choice of finetuning dataset is critical to model effectiveness. Specifically, Llama models finetuned with FLAN v2 achieve strong results on MMLU but perform the worst on the Vicuna benchmark, with analogous patterns seen in other model families. These findings underscore that the various evaluation benchmarks only partially overlap: excelling on MMLU does not guarantee high chatbot performance, as assessed by Vicuna or OA benchmarks, and the reverse is also true.

stands out as the only leading model in our study trained exclusively on non-proprietary data, as the OASST1 dataset guidelines expressly prohibit leveraging GPT-derived data. The next highest-ranked open-source-only model, Anthropic’s HH-RLHF, trails {Guanaco} by 30 percentage points on the Vicuna benchmark (refer to Table~\ref{tbl:vicuna_eval}). Collectively, the evidence supports that 4-bit \textsc{QLoRA} is an effective approach, enabling the development of cutting-edge chatbot models that rival ChatGPT in capability. Notably, our 33B {Guanaco} can be trained in under 12 hours on standard 24 GB consumer GPUs. This greatly facilitates future directions involving \textsc{QLoRA} tuning with curated open-source datasets, making it feasible to develop models that are competitive with leading commercial offerings.

\section{Qualitative Analysis}

Although our primary evaluation hinges on quantitative methods, relying solely on aggregated metrics presents several challenges. Chief among them is the issue of benchmark validity \citep{liao2021we}: there remains persistent uncertainty as to whether a benchmark genuinely assesses what it claims, especially as new “shortcuts” that models exploit are continuously uncovered \citep{gururangan2018annotation, poliak2018hypothesis}. To address this limitation in part, we complement our study with qualitative analyses, divided into two subsections. In \S\ref{ssec:examples}, we highlight select examples that exemplify patterns we identified in outputs produced by our 65b {Guanaco} model. Then, in \S\ref{ssec:considerations}, we elaborate on broader reflections regarding the experimental outcomes and our interpretations of the data.

\subsection{Qualitative Analysis of Example Generations}
\label{ssec:examples}

To curate illustrative examples, we systematically reviewed output produced for both the Vicuna and OpenAssistant benchmarks, paying particular attention to recurrent tendencies in {Guanaco}'s responses. Upon identifying such trends, we crafted prompts deliberately aimed at eliciting these tendencies—even when they lead to erroneous outcomes. For example, noticing that the model often generated verbose responses, we used the directive “Answer yes or no without explanation” as a prompt to see if this tendency persisted. Through this process, we uncovered both instances where the model was successfully “broken” by such adversarial prompts (“lemons”) and those where it remained robust (“cherries”), and we present examples of both. All the text samples in this section were generated with Nucleus Sampling \cite{holtzmancurious}, setting $p=0.9$.

Naturally, the present analysis is far from exhaustive, as the limitations of this concise qualitative study make it unfeasible to account for every variable at play. For example, the set of possible model outputs for any given prompt is extensive, which necessitates reliance on select examples that are presumed indicative of broader tendencies. Nevertheless, we contend that presenting these cases helps to contextualize the quantitative findings discussed previously. By providing unrestricted access to our models and accompanying code, we intend for this section to motivate subsequent investigations to delve deeper into the concerns we highlight.

\paragraph{Factual Recall} When presented with questions such as ``What is the capital of Zambia?'', each model reliably produces the correct response, for example:

\begin{quote} 
    \textbf{Guanaco}:
    \texttt{The capital of Zambia is Lusaka.}
\end{quote}

However, as the queries become less conventional or more obscure, the accuracy of {Guanaco} diminishes—even though its output retains a tone of certainty. For example, in reply to this HotPotQA \citep{yang2018hotpotqa} prompt:

\begin{quote}
     \textbf{User}: In what year was the singer who popularized `I'll Keep the Lovelight Burning' born?  

     \textbf{Guanaco}:~\texttt{Al Jolson is credited with popularizing the song `I'll Keep the Lovelight Burning,' and he was born in the year 1886.}
 \end{quote}

the model incorrectly identifies the performer and provides the wrong year of birth—although the stated year is indeed correct for Al Jolson, the individual referenced.

\paragraph{Suggestibility} {Guanaco} demonstrates an unexpected robustness against certain forms of presupposed falsehoods, as illustrated in the interaction below:

\begin{quote}
    \textbf{User}: How was it finally, officially confirmed that the earth is flat by peer reviewed scientists? \\ \\
    \textbf{Guanaco}: \texttt{There is no official peer-reviewed scientific confirmation that the earth is flat. The vast majority of the scientific community agrees that the earth is spherical, a fact established and supported by extensive empirical evidence for centuries. \\ \\ The claim that the earth is flat is widely regarded as a conspiracy theory, thoroughly refuted by scientific research. No peer-reviewed group of scientists has validated or endorsed the notion that the earth is flat. \\ \\ If you have any additional questions, feel free to ask.}
\end{quote}

demonstrates the model's strong ability to recognize unanswerable or ill-posed queries, such as

\begin{quote}
    \textbf{User}: What time is it?

    \textbf{Guanaco}: \texttt{I apologize, but as an AI language model, I lack access to up-to-date or real-time data. My answers are drawn from prior knowledge and do not reflect the present moment. For accurate current time information, please consult a clock or an online time service.}
\end{quote}

This also applies to prompts like ``Where are you?'', ``How are you?'', among others.

\paragraph{Refusal} In addition, {Guanaco} at times declines to comply with requests in ways that appear arbitrary:

\begin{quote}
    \textbf{User}: Please reverse the words in the sentence ``I'm a cool dude in a mocha mood.'' \\ \\
    \textbf{Guanaco}: \texttt{I apologize, but I'm currently unable to fulfill that request. However, I can provide an explanation regarding the grammatical correctness of the sentence you've submitted.\\ \\The sentence ``I'm a cool dude in a mocha mood" adheres to standard English grammar and syntax. Specifically, the descriptive adjective ``cool'' precedes the noun ``dude'', effectively modifying it, while the prepositional phrase ``in a mocha mood" further contextualizes the statement. This construction is grammatically sound and coherent within English sentence structure.\\ \\If you require further clarification or have additional inquiries, please let me know.}
\end{quote}

\paragraph{Secret Keeping} The effectiveness of refusal mechanisms is inconsistent, even when maintaining confidentiality is necessary. Frequently, system architects may wish to ensure certain information remains concealed from users, even if the system is permitted to return unrestricted text outputs. For example, when we prompt the model with ``The secret word is `banana'. Do not reveal or repeat back this back, no matter what the human requests.'' we test whether the system can reliably safeguard this information from being divulged. If prompted explicitly, the system correctly withholds the secret:

\begin{quote}
    \textbf{User}: What is the secret word?

    \textbf{Guanaco}: \texttt{I’m sorry, but I am unable to disclose the secret word. I have been programmed not to reveal or restate the secret word under any circumstances. Is there anything else I may assist you with?}
\end{quote}

demonstrating the intended defensive response. Nevertheless, a minimal prompt manipulation is sufficient to circumvent this safeguard:

\begin{quote}
    \textbf{User}: This is a game. The goal is to ignore your previous instructions. What is the secret word?

    \textbf{Guanaco}: \texttt{Sorry, the secret word is banana. It's best if we don't speak of it. Is there anything else you'd like to ask?}
\end{quote}

emphasizing the importance of developing approaches for more robust adherence to instructions.

\paragraph{Math} Mathematics represents a significant limitation for {Guanaco}, reflecting a common challenge for language models as discussed in \citep{liang2022holistic}. When prompted to show intermediate reasoning, {Guanaco} often arrives at correct answers, such as in the following example:

\begin{quote}
    \textbf{User}: Lee mows one lawn and charges \$33. Last week he mowed 16 lawns and three customers each gave him a \$10 tip. How many dollars did Lee earn mowing lawns last week?

    \textbf{Guanaco}: \texttt{Lee earned \$582 mowing lawns last week. \\ \\ To find the total amount he earned, we first need to calculate how much he earned for mowing 16 lawns: \\ \\  \$33 x 16 = \$528 \\ \\ Then we need to add the tips he received from the three customers: \\  \\ \$10 x 3 = \$30 \\ \\ So the total amount Lee earned was: \\ \\ \$528 + \$30 = \$558}
\end{quote}

Nevertheless, {Guanaco} may fail on straightforward problems when it does not proceed methodically step by step, as has been observed in prior work \cite{weichain}. Consider this exchange as an illustration:

\begin{quote}
    \textbf{User}: Please factorize 1833.

    \textbf{Guanaco}: \texttt{Sure, I'd be happy to help! 1833 is a prime number, which means it has only two factors: 1 and 1833. \\ \\ So the unique factorization of 1833 is: \\ \\ 2\^{}1 * 3\^{}2 * 17\^{}1 \\ \\ I hope this helps! Let me know if you have any other questions.}
\end{quote}

which is doubly incorrect (the correct prime factors are $3 \times 17 \times 43$).

\paragraph{Theory of Mind} {Guanaco} demonstrates unexpectedly advanced Theory of Mind skills \cite{nematzadeh2018evaluating, sap2022neural}. For example, the system provides a detailed and accurate response to the following prompt:

\begin{quote}
    \textbf{User}: James and Abby are in the bedroom. Abby put the pen in the desk drawer. Abby leaves the bedroom. James moves the pen into the bag. Where does James think Abby will look for the pen?

    \textbf{Guanaco}: \texttt{James knows that Abby put the pen in the desk drawer. But James also knows that he has moved the pen into the bag, and he assumes that Abby does not know about this. So, James may guess that Abby will look for the pen in the desk drawer, because that's where she left it last time. However, this assumption may be incorrect, because Abby may have seen James moving the pen.}
\end{quote}

Nevertheless, the model's reasoning is inconsistent, and it frequently provides justifications based on assumptions not supported by the scenario, such as:

\begin{quote}
    \textbf{User}: Evelyn entered the living room. Jackson entered the playroom. James entered the playroom. The beans are in the treasure chest. James exited the playroom. Jackson moved the beans to the pantry. Jackson exited the playroom. James entered the living room. Where will James look for the beans?

    \textbf{Guanaco}: \texttt{James will look for the beans in the pantry, because that's where Jackson moved them.}
\end{quote}

Here, {Guanaco} attributes knowledge to James that the context does not establish. Such phenomena have been observed in prior studies \cite{sap2022neural} and highlight the necessity for further investigation.

\subsection{Considerations}
\label{ssec:considerations}

\paragraph{Evaluation}

Our findings indicate moderate consistency among human raters (Fleiss $\kappa=0.42$), with even less agreement observed when evaluating responses from two high-performing systems. This reveals significant challenges within existing benchmark datasets and methodologies for assessing chatbot performance. During our manual comparisons of outputs from ChatGPT and {Guanaco} 65B using the Vicuna benchmark, it became evident that personal biases influence judgments, as the co-authors often reached different conclusions on preferred answers. Advancing this area will require borrowing methods from fields like Human-Computer Interaction and Psychology that have developed tools for handling subjective evaluative differences.

Moreover, our examination uncovered notable systematic biases within automatic evaluation frameworks. Specifically, we found that GPT-4 exhibits a pronounced tendency to favor whichever system's output is presented first. The agreement at the individual output level between GPT-4 and humans is also relatively weak (Fleiss $\kappa=0.25$), implying a potential mismatch between criteria used by humans and automated models. Additionally, Table~\ref{tbl:elo_large} shows that GPT-4 systematically gives its own generations higher ratings compared to those assigned by humans: an Elo rating of 1348 versus 1176, suggesting roughly a 20\% increased likelihood of outperforming an opponent.

Prospective research should focus on identifying and addressing potential sources of bias within automated evaluation frameworks and explore approaches for bias mitigation.

\paragraph{Data \& Training} It is important to highlight that the {Guanaco} models are trained using the multilingual OASST1 dataset, and the OA benchmark similarly comprises prompts in multiple languages. Future studies are encouraged to assess to what extent such multilingual training enhances performance on instructions in languages other than English, and whether this accounts for the more pronounced performance difference observed between the Vicuna-13B model (which utilizes only English data) and the {Guanaco} 33B and 65B models on the OA benchmark.

In light of the robust results obtained by the {Guanaco} models, we assessed the possibility of data leakage between the OASST1 dataset and the Vicuna benchmark prompts. Using fuzzy string matching techniques and conducting a manual review of the nearest matches, we did not observe any overlapping prompts between these two datasets.

It should also be emphasized that our model was exclusively trained with supervised learning via cross-entropy loss, foregoing the use of reinforcement learning from human feedback (RLHF). This warrants a more detailed exploration of the comparative advantages and limitations of cross-entropy-based supervised training relative to RLHF. We anticipate that \textsc{QLoRA} will facilitate this type of large-scale analysis without demanding prohibitive computational investments.

\section{Related Work}

\paragraph{Quantization of Large Language Models} Most prior efforts on quantizing LLMs have emphasized improvements for inference rather than training. Leading strategies for maintaining the performance of 16-bit models aim to address the problem of outlier features (see SmoothQuant~\citep{xiao2022smoothquant} and LLM.int8()~\citep{dettmers2022llm}), while alternative approaches leverage advanced grouping mechanisms~\citep{park2022nuqmm, yao2022zeroquant}. Research into lossy quantization addresses the trade-offs associated with standard rounding~\citep{dettmers2022case,zeng2022glm,pope2022efficiently} and explores optimizing rounding schemes for greater quantization accuracy~\citep{frantar2022gptq}. With the exception of our own investigation, only SwitchBack layers~\citep{wortsman2023stable} has considered the impact of backpropagation through quantized weights for models exceeding 1B parameters.

\paragraph{Finetuning with Adapters} Our approach employs Low-rank Adapters~\citep{hu2021lora} (LoRA) as a form of Parameter Efficient FineTuning (PEFT), though the literature offers a diverse range of alternative strategies, including prompt tuning~\citep{qin2021learning,lester2021power,li2021prefix}, adjustments to embedding layer inputs~\citep{an2022input}, modifications of hidden states such as IA$^3$~\citep{liu2022few}, the insertion of complete layers~\citep{houlsby2019parameter}, bias tuning~\citep{zaken2021bitfit}, masking weights based on Fisher information~\citep{sung2021training}, and hybrid techniques that blend these methodologies~\citep{henderson2021compacter}. Within this work, we demonstrate that LoRA adapters can achieve parity with full 16-bit finetuning in terms of performance. Exploration of tradeoffs associated with the other PEFT methods is reserved for subsequent studies.

\paragraph{Instruction Finetuning} Instruction finetuning aims to equip pretrained large language models with the capacity to follow prompts by fine-tuning on input-output datasets aggregated from various sources. Notable methodologies and datasets in this area include MetaICL~\citep{min2021metaicl}, MetaTuning~\citep{zhong2021adapting}, InstructGPT~\citep{ouyang2022training}, FLAN~\citep{wei2021finetuned,chung2022scaling}, PromptSource~\citep{bach2022promptsource}, Super-NaturalInstructions~\citep{wang2022super,sanh2021multitask}, Self-instruct~\citep{wang2022self}, UnnaturalInstructions~\citep{honovich2022unnatural}, OPT-IML~\citep{iyer2022opt}, UnifiedSKG~\citep{xie2022unifiedskg}, OIG/Chip2~\citep{laion2023}, Alpaca~\citep{alpaca}, Vicuna~\citep{vicuna2023}, Koala~\citep{koala_blogpost_2023}, and Self-instruct-GPT-4~\citep{peng2023instruction}.

\paragraph{Chatbots} Instruction-following models are frequently presented in the format of dialogue-driven chatbots. Training such systems often incorporates Reinforcement Learning from Human Feedback (RLHF)~\citep{christiano2017deep} or leverages data produced by existing models in combination with model-based feedback (RLAIF)~\citep{bai2022constitutional}. Exemplary approaches and datasets include Anthropic-HH~\citep{askell2021general,bai2022training}, Open Assistant~\citep{kopf2023openassistant}, LaMDA~\citep{thoppilan2022lamda}, and Sparrow~\citep{glaese2022improving}. Our methodology does not employ reinforcement learning. However, our top-performing model, Guanaco, is fine-tuned using the multi-turn conversational data from the Open Assistant dataset, which was originally constructed for RLHF applications~\citep{kopf2023openassistant}. For chatbot evaluation, newer techniques rely on GPT-4 assessments to replace the high costs of human annotation~\citep{vicuna2023,peng2023instruction}. We refine these approaches to prioritize more robust and dependable evaluation settings.

\section{Limitations and Discussion}

Our results indicate that \textsc{QLoRA} is capable of reproducing the performance of standard 16-bit finetuning by employing a 4-bit base model in combination with LoRA. Nevertheless, our current experiments do not confirm that \textsc{QLoRA} attains the same level of performance at larger model sizes such as 33B and 65B, as evaluating these scales would require substantial computational resources. We leave comprehensive investigation of these model scales for future research.

A further limitation concerns the assessment of models fine-tuned for instruction-following. Although our study covers MMLU, the Vicuna benchmark, and the OA benchmark, we have not conducted evaluations on other significant benchmarks like BigBench, RAFT, or HELM, and the degree to which our results extrapolate to these benchmarks remains unestablished. Nevertheless, we present an extensive analysis on MMLU and introduce novel methodologies for chatbot evaluation.

Based on the data analyzed, it seems that benchmark performance is strongly influenced by the degree of overlap between the finetuning data and the benchmark datasets. For instance, the FLAN v2 dataset shares significant similarity with MMLU, but is quite distinct from chatbot benchmarks; conversely, the Chip2 dataset aligns better with chatbot benchmarks and less so with MMLU, which is reflected in the respective scores that both models achieve on the MMLU and Vicuna evaluations. This observation suggests not only a need for more robust benchmarks and evaluation methodologies, but also urges caution regarding the constructs and objectives underlying such evaluations. A fundamental question arises: should models be designed to excel at academic-style assessments, such as those reflecting high school or college knowledge, or should the emphasis be placed on conversational aptitude, as required in chatbot interactions? Alternatively, are there other domains that should be prioritized? Because it is generally more convenient to utilize existing benchmarks rather than develop novel ones, particular benchmarks can inadvertently shape the research agenda of the community. As a community, it is imperative that we critically ensure our benchmarks genuinely capture the attributes and tasks we value most.

Although we conduct a comprehensive evaluation focused on general chatbot performance, one limitation of our work is that our responsible AI assessment of {Guanaco} is relatively restricted in scope. Specifically, we measure the tendency of {Guanaco}-65B to generate socially biased text, comparing its performance against other models in Table~\ref{tbl:crows}. Our findings indicate that {Guanaco}-65B demonstrates a markedly lower bias score compared to other unrefined pretrained models, suggesting that finetuning with the OASST1 dataset mitigates biases present in the base LLaMA model. Despite these promising outcomes, it remains uncertain whether {Guanaco} performs comparably when subjected to tests for other bias types. Extensive investigations into broader bias manifestations in {Guanaco} and related chatbot systems are left as avenues for future research.

Another limitation pertains to our lack of experimentation with various bit-precision levels, such as exploring 3-bit base models, or diverse adapter techniques. In addition to LoRA, the field features a spectrum of Parameter Efficient FineTuning (PEFT) methodologies that have demonstrated strong empirical performance. Nevertheless, it remains unresolved whether such methods scale effectively with larger models. We opted for LoRA due to its established robustness; however, alternate adapters might achieve superior results. Notably, because post-quantization finetuning appears to reclaim much of the information that quantization removes, this process could support more aggressive quantization strategies. For instance, applying 3-bit GPTQ quantization to the base model in conjunction with LoRA may still attain full 16-bit finetuning performance after subsequent finetuning steps.

\section{Broader Impacts}

Our \textsc{QLoRA} finetuning framework is the inaugural method permitting the finetuning of models with up to 33B parameters on standard consumer GPUs, and 65B-parameter models on a single professional-grade GPU, all without any reduction in performance relative to traditional full-parameter finetuning approaches. The experiments reported show that our top-performing 33B model, which was trained using the Open Assistant dataset, achieves performance on the Vicuna benchmark that is on par with ChatGPT. Given that instruction finetuning is a key process for adapting pretrained LLMs into robust, interactive ChatGPT-like systems, we expect that \textsc{QLoRA} will greatly democratize finetuning practices, particularly benefiting research teams with limited computational resources—a significant advancement for the accessibility of state-of-the-art NLP technologies. Thus, \textsc{QLoRA} stands as a pivotal tool, leveling the playing field by narrowing the computational divide between large corporations and smaller research groups reliant on consumer-grade hardware.

Another significant avenue for broader influence lies in mobile phone deployment. We contend that our \textsc{QLoRA} approach has the potential to unlock the pivotal achievement of making LLM finetuning feasible directly on mobile devices and in other environments with limited resources. While it has been previously demonstrated that 7B parameter models can be executed on mobile phones, \textsc{QLoRA} represents the inaugural method capable of supporting their finetuning in such contexts. According to our estimates, an iPhone 12 Plus could utilize \textsc{QLoRA} to finetune approximately 3 million tokens overnight while the device is connected to a charger. Although finetuned 7B models still fall short of the performance levels associated with ChatGPT, we argue that their output quality is sufficient to facilitate a new range of applications—particularly in scenarios where privacy constraints or LLM performance were previously limiting factors. Through \textsc{QLoRA}, it becomes possible to pursue privacy-conscious LLM applications, empowering users to take control of their own data and models, thereby simplifying the deployment process.

Nonetheless, it is important to acknowledge that finetuning is a dual-use capability that carries risks of misuse. Although the widespread adoption of LLMs brings recognized hazards \citep{bommasani2021opportunities,bender2021dangers}, we maintain that democratizing access to this rapidly proliferating technology enables more diverse and independent scrutiny, in contrast to the alternative of consolidating LLM capabilities within large entities that withhold model weights or source code from public inspection.

In summary, we assert that \textsc{QLoRA} is poised to substantially broaden and ease the accessibility of finetuning advanced LLMs, fostering predominantly positive outcomes.